\chapter{The Harness: The Layer the Contest Moved To}
\chaptermark{The Harness}
\label{ch:harness}

Every chapter so far has treated the model as the object of competition. This chapter argues that the object is moving, and that the move is the most strategically consequential development of 2026---more consequential than any individual release, because it changes what winning a release would even mean.

The claim in one sentence: \emph{models are swappable; harnesses are programmed to}. A model sits behind an API call and can be exchanged for another with a configuration change---an entire product category, the LLM gateway, exists to make exactly that substitution trivial. A \emph{harness}---the runtime that holds the agent loop, the tool schemas, the permission system, the memory and context management, the sandbox, the session store, the subagent orchestration---is something an organization writes code \emph{against}. Its abstractions become the shape of the integration. Its execution semantics become the assumptions the tests encode. Its permission model becomes the artifact an auditor reads. If that is right, then the layer where switching is expensive is not the layer everyone is measuring, and a strategy that wins the harness can lose every benchmark and still take the ecosystem.

\section{The layer acquires a name}

A useful sign that an industry has discovered a layer is that its competitors independently start using the same word for it. Through 2026 all three of the leading agentic laboratories published engineering material on \emph{harnesses} as such: OpenAI on harness engineering in an agent-first world, Anthropic on harness design for long-running application development, and---in August---DeepSeek, in a paper that cites both~\cite{cordis-paper}. Three laboratories that agree on almost nothing agree that this is a distinct artifact with its own design discipline [V].

What is inside it is mostly not the model call. The most careful public anatomy of a production harness, an analysis of Claude Code by researchers at MBZUAI and UCL, puts the finding plainly: ``the core of the system is a simple while-loop that calls the model, runs tools, and repeats. \emph{Most of the code, however, lives in the systems around this loop}: a permission system with seven modes and an ML-based classifier, a five-layer compaction pipeline for context management, four extensibility mechanisms (MCP, plugins, skills, and hooks), a subagent delegation and orchestration mechanism, and append-oriented session storage''~\cite{claude-code-anatomy}. A widely repeated corollary---that only 1.6\% of the codebase is AI decision logic---should be handled with more care than it usually receives: the paper cites it as a \emph{community estimate} derived from line-counting a leak-derived bundle containing generated and minified code, and does not present it as a measurement [A]. The qualitative claim survives the quantitative one, and the qualitative claim is enough.

\section{The harness is a larger economic lever than the model}

The strongest evidence that this layer matters is not rhetorical but experimental, and as of this edition there is exactly one controlled study of it. Researchers at Writer held twenty-two audited evaluation tasks, six foundation models from five vendors across three capability tiers, and identical judges, prices, and prompts constant, and varied only the orchestration layer: a baseline production loop against a purpose-built harness~\cite{harness-effect}. Blended across models, the harness cut cost per task by 41\%, tokens per task by 38\%, and median latency by 44\%, while raising quality per dollar by 82\%, at maintained task completion.

The sentence that matters for this book is the comparison the authors themselves draw: \emph{keeping any model and adopting the harness saved between 33\% and 61\%, while switching between the cheapest and most expensive model under baseline conditions saved only 36\%} [V, single study]. Every model in the panel got substantially cheaper regardless of vendor or tier. If that result generalizes---and twenty-two tasks from the vendor of the winning harness is a caveat that must be printed beside it---then the harness is a \emph{larger} lever on production economics than the choice of model, which inverts the entire public conversation about which layer is competitive.

There is a second, quieter consequence. A harness is also the instrument by which models are measured: the same model produces different results under different harnesses, and system-prompt budgets across harnesses vary by nearly two orders of magnitude~\cite{harness-register}. Whoever owns the harness owns the conditions of measurement---which is worth remembering the next time a scoreboard is cited, including in this book.

There is a third consequence, and it compounds rather than competes with the first two. Chapter~\ref{ch:compute}'s priced bill of materials (\S\ref{sec:appliance-cost}) shows a capacity-first local appliance running the largest open-weight model in existence at a little over a dollar per million output tokens, against twenty-five to thirty for the closed frontier---and, at realistic utilization, below a hyperscaler's own marginal cost of serving the identical model. If that holds, model cost is heading toward a rounding error at the margin for an organization willing to run its own hardware, which does not shrink the harness argument, it \emph{sharpens} it: the Writer study's 33--61\% saving from adopting a harness while holding the model constant was already larger than the 36\% saving from switching models under a fixed harness; once the model's own marginal cost approaches zero, the harness is not merely the larger of two levers, it is closer to the \emph{only} lever left on production economics. An organization that owns both the cheap inference tier and the harness that drives it captures the saving twice, which is the strongest economic reason---stronger than the licence, stronger than the star count---to expect DeepSeek Harness and its peers (Table~\ref{tab:harness}) to be paired with exactly the kind of local capacity Chapter~\ref{ch:compute} prices, rather than deployed against a metered API.

\section{What China shipped in August}

On 13 August 2026 DeepSeek released \emph{DeepSeek Harness}, an MIT-licensed agent harness whose organizing principle is that \emph{everything is a plugin}: models, tools, skills, sessions, sandboxes, filesystems, loops, orchestration, and interface are all components, with no privileged core to patch~\cite{dsh-repo,dsh-official}. It is model-agnostic by construction---Anthropic, OpenAI, Bedrock, Azure, Gemini and DeepSeek endpoints are all supported---it speaks MCP through a client plugin (tools only; resources and prompts are not bridged), and it ships subagent bridges that let it drive Claude Code and OpenAI Codex as subordinate components~\cite{dsh-mcp,dsh-newstack}. Within roughly forty-eight hours the repository passed 100{,}000 GitHub stars [V]. For scale: LangGraph, the most-downloaded open agent framework, has accumulated roughly 34{,}000 stars over more than a year.

The accompanying paper is stranger and more interesting than a release note. \emph{A Programming Paradigm for Spatiotemporal Composability}, an eighty-eight-page draft from Peking University and DeepSeek, is a programming-languages paper, not a machine-learning one~\cite{cordis-paper}. Its subject is a meta-framework called Cordis, and its two contributions are \emph{temporal composability}---unloading a component provably returns the system to the state it would have occupied had the component never loaded, via \emph{revertible effects} in which every context mutation carries a runtime-tracked inverse---and \emph{spatial composability}, in which components declare dependencies and are reactively notified as those dependencies change, via \emph{reactive coeffects}. The lineage is Moggi on monads, Plotkin and Power on algebraic effects, Petricek and Orchard on coeffects; the novelty claim is lifting these static type-system ideas into runtime mechanisms. Two honest qualifications belong here, both stated by the authors: the paper is a self-hosted preprint ``under active revision'' rather than a refereed publication, and its validation case study is observational rather than controlled. It is also worth being precise that the paper describes \emph{Cordis}, the substrate; the harness is a product built on it, and applying the formalism to self-evolving harnesses is named as future work, not as an achieved result.

The provenance is where the strategic content sits. Cordis is not new and did not originate at DeepSeek: it is the kernel of Koishi, a four-year-old Chinese open-source chatbot framework with over four thousand community-contributed plugins, and the paper's first author is Koishi's creator~\cite{cordis-paper,dsh-36kr}. DeepSeek did not build a plugin ecosystem; it recruited the maintainer of one, formalized his runtime with a university partner, and shipped it under MIT. That is an ecosystem-acquisition strategy rather than a model strategy, and it is this book's thesis executed one layer up.

The complements logic is exact, and it is the mechanism of Chapter~\ref{ch:theory} rather than a new one. DeepSeek gave the harness away under the most permissive licence available \emph{on the same day} it raised API prices for V4-Pro and V4-Flash and introduced peak/off-peak pricing~\cite{dsh-venturebeat,deepseek-v4pro}. Commoditize the complement, monetize the scarce layer, and---because the harness is model-agnostic---ensure that the commoditization damages competitors' pricing power at least as much as your own. The \emph{South China Morning Post} read the move as ``shifting the global competition from model intelligence to one that focuses on how seamlessly an AI agent can connect to real-world software''~\cite{dsh-scmp}, which is the thesis of this chapter stated by a newspaper.

DeepSeek is not alone, and the inventory matters more than any single entry (Table~\ref{tab:harness}). What the table shows is a bloc shipping runtimes, sandboxes, command-line harnesses, and agent platforms---infrastructure rather than weights---across at least four laboratories and one state initiative.

\begin{table}[htbp]
  \centering
  \small
  \caption{Chinese agent-infrastructure artifacts as of \datafreeze. The pattern is runtimes and protocols, not only models.}
  \label{tab:harness}
  \begin{tabular}{p{34mm}p{22mm}p{38mm}p{22mm}p{16mm}}
    \toprule
    Artifact & Organization & Type & Licence & Grade \\
    \midrule
    DeepSeek Harness (\texttt{dsh}) & DeepSeek & agent harness, plugin-native & MIT & V \\
    Cordis & DeepSeek / PKU & meta-framework + formal paper & open & V \\
    OpenSandbox & Alibaba $\to$ neutral org & sandbox runtime for agents & Apache 2.0 & V \\
    Qwen Code, Qoder CLI & Alibaba & CLI coding harnesses & Apache 2.0 & A \\
    Qianwen Open Platform & Alibaba & third-party agent platform & proprietary & V \\
    AgentLoop / Teams / Run & Alibaba Cloud & agent operations stack & proprietary & V \\
    Kimi Code CLI & Moonshot & CLI harness, ACP-compatible & MIT & A \\
    CodeBuddy Code & Tencent & CLI harness, ACP-compatible & --- & A \\
    Dify & LangGenius & agent/RAG platform, 144k stars & open & A \\
    \midrule
    Agent interoperability initiative & CAC (state) & standards-diplomacy instrument & --- & V \\
    \bottomrule
  \end{tabular}
\end{table}

\section{The standards layer, and who is not in the room}

The governance picture is the sharpest geopolitical fact in this chapter, and it is a fact about membership lists. Anthropic donated the Model Context Protocol to the Linux Foundation's newly created Agentic AI Foundation in December 2025, co-founded with Block and OpenAI, with Google, Microsoft, AWS, Cloudflare and Bloomberg as supporting members; MCP at donation reported more than ten thousand active public servers and over ninety-seven million monthly SDK downloads~\cite{mcp-aaif}. Google's Agent2Agent protocol, donated to the Linux Foundation, passed 150 organizations in its first year~\cite{a2a-year}. In August 2026 a further standard, Agent Plugins 1.0, arrived with a technical steering committee drawn from Amazon, Cursor, Microsoft, OpenAI, Vercel and Google~\cite{agent-plugins}.

\emph{No Chinese company is a named member of any of them} [V, by absence]. And on 17 July 2026 the Cyberspace Administration of China released a ``Global Cooperation Initiative on Agent Mutual Trust, Interconnection, and Interoperability,'' calling for ``open, non-discriminatory, and transparent international interoperability standards for agents'' and for shared agent identity and interface-compatibility infrastructure---while naming no existing protocol, neither MCP nor A2A~\cite{cac-agents}. Two weeks later DeepSeek shipped an MIT harness that reached a hundred thousand stars in two days.

Chapter~\ref{ch:governance} argued that standards are where governance becomes irreversible, and gave China's standards strategy as the arena the West systematically under-weights. This is that argument arriving at the agent layer, and the structure is the familiar one from Part~I: the incumbent bloc convenes a consortium of incumbents; the challenger declines to join, ships a permissively licensed implementation, and bids for the installed base directly. Whether that succeeds is unknown. That it is the same play, run again, on a layer nobody was watching a year ago, is not.

\section{The strongest case against this chapter}

This book's discipline requires the counter-argument at full strength, and here it is unusually strong---strong enough that the thesis must be narrowed rather than merely defended.

\emph{First, the interface is actively being commoditized, right now, by the incumbents.} Agent Plugins 1.0---announced nine days before this edition---is a vendor-neutral package format for distributing skills and MCP servers across agent clients ``without requiring file rearrangement or manifest rewriting,'' governed independently and developed in public~\cite{agent-plugins}. Six of the largest firms in the industry are standardizing precisely the interface this chapter calls sticky. A thesis that harness lock-in is structural has to explain why the participants are dismantling it.

\emph{Second, the architecture is not a secret.} Four independent teams---Claude Code in TypeScript, Codex CLI in Rust, Aider in Python, and OpenClaw---converged on the same structural pattern: an outer loop, tool execution, result capture, staged permissions~\cite{claude-code-anatomy}. Convergent design is evidence that the architecture is obvious, and obvious architectures are not moats.

\emph{Third, the migration evidence points to weeks, not quarters.} The best-documented public account of leaving a framework---Octomind's departure from LangChain---reports rigid abstractions, debugging overhead, and a loss of iteration speed, but quantifies no switching cost~\cite{octomind}. More pointedly, OpenAI's Codex shipped a one-line importer for Claude Code configuration: the cost was real enough to automate, and small enough to be automated in one command~\cite{harness-register}. Skills appear to transfer better than the thesis wants: one reported comparison found a skill trained under one harness scoring 80.4 when transferred versus 81.8 trained natively [A, one skill, one benchmark pair].

\emph{Fourth, the builder of the most successful harness in the world denies the premise.} Boris Cherny, the creator of Claude Code, has described the harness as ``the thinnest possible wrapper'' and located ``all the secret sauce'' in the model~\cite{harness-register}. Against him, an engineer at Obvix Labs supplies the opposing epigram---``the model is interchangeable, the harness is not.'' Both cannot be right, and the person with the most direct evidence is on the other side of this chapter's argument.

\emph{Fifth, and most awkward for a chapter about Chinese lock-in strategy: Alibaba is doing the opposite.} OpenSandbox---the agent execution sandbox released in March 2026, not August---was moved out of Alibaba's GitHub organization into neutral governance, carries an OpenSSF best-practices badge, is listed on the CNCF landscape, and lists Claude Code, Gemini CLI, LangGraph and Google ADK among its supported integrations~\cite{opensandbox}. That is a firm commoditizing the execution layer beneath every harness, American ones included, and donating it away---the Kubernetes playbook, not a capture play. Any account that treats the Chinese move as uniformly lock-in-seeking is contradicted by its largest participant.

\emph{Sixth, and decisively for intellectual honesty: nobody has measured the thing this chapter asserts.} There is no published, methodologically sound study comparing the cost of switching harnesses against the cost of switching models. The Writer experiment measures operating economics, not migration cost. Enterprise surveys establish that multi-model deployment is now the norm---over three-quarters of teams run multiple models---but do not measure framework-switching rates at all~\cite{langchain-survey}. This chapter therefore advances a \emph{well-motivated hypothesis with strong circumstantial support}, and it should be read at that grade [S], not as a finding.

\section{The thesis, narrowed to what survives}

What survives the counter-argument is not the strong claim but a sharper one. The technical interface \emph{is} commoditizing, quickly, and MCP plus Agent Plugins will likely make tool definitions and skill packages portable across clients. Lock-in does not therefore disappear; it \emph{relocates}---which is the rent-relocation mechanism of Section~\ref{sec:io} applied to infrastructure rather than to weights. It relocates to three places the standards do not touch.

\emph{To distribution.} Agent Plugins standardizes packaging, not distribution. Each platform still controls its own marketplace, so an author of a valuable skill is not merely building to a standard---they are choosing which gatekeeper reaches enterprise buyers. The switching cost becomes commercial rather than technical, which is harder to automate away than a config importer. As one practitioner put it, if useful behavior ends up inside vendor namespaces, ``the standard will look open while lock-in just moves elsewhere''~\cite{agent-plugins-newstack}.

\emph{To the evaluation substrate.} If the harness determines measured model quality, then the harness that an ecosystem standardizes on defines the scoreboard that ecosystem believes. DeepSeek used its harness internally to benchmark its own coding models before releasing it [A]. Owning the instrument is a durable position that no packaging standard addresses.

\emph{To certification, which is where this argument meets Chapter~\ref{ch:trust}.} For unregulated development, the transfer evidence suggests re-validation is cheap. For regulated deployment it is not, and for a reason that has nothing to do with performance regression: the artifact an auditor examines is the permission model, the sandbox boundary, the audit log, and the tool-invocation trace---all of which are harness properties, not model properties. Re-certifying an agent system after a harness migration is paperwork, evidence, and elapsed calendar time at an assessor, and the certification ladder of Section~\ref{sec:certladder} prices none of it against the model. This is the version of the switching-cost argument that most plausibly holds at scale, and it is testable: if harness migrations in regulated sectors run materially longer than in unregulated ones, the mechanism is real.

Three implications follow for this book's propositions. P1 (adoption) acquires a second front: adoption of Chinese \emph{infrastructure} is now measurable independently of adoption of Chinese weights, and the star counts say it is moving faster. P3 (rent destruction) is reinforced---a hundred-thousand-star MIT harness released the same day as an API price increase is the rent-relocation thesis performed rather than argued. And the book's central asymmetry sharpens into its most uncomfortable form: \emph{a bloc can lose the model contest and still win the ecosystem}, because the thing an enterprise cannot cheaply change is not which model answers the call but which runtime made the call, logged it, sandboxed it, and can prove to an auditor what it did. If American models remain the best---and on today's independent indices they are---and the world's agents are orchestrated by runtimes shipped from Hangzhou and Beijing under MIT licences, then superiority at the model layer buys less than its owners expect. That is the claim this chapter registers, at scenario grade, with its falsifiers stated in Appendix~\ref{app:dashboard}.
